% paper.tex
%
% Workshop submission for ICML 2026 (Mechanistic Interpretability track).
% Page limit: 4--5 pages excluding references.
%
% Drop this file alongside the ICML 2026 LaTeX style files (icml2026.sty,
% icml2026.bst, fancyhdr.sty, etc.) and compile with pdflatex + bibtex.

\documentclass{article}
\usepackage{microtype}
\usepackage{graphicx}
\usepackage{subfigure}
\usepackage{booktabs}
\usepackage{hyperref}
\usepackage{url}
\usepackage{amsmath}
\usepackage{amssymb}
\usepackage{algorithm}
\usepackage{algorithmic}
\usepackage{xcolor}
\usepackage{multirow}

% Use the ICML 2026 style (anonymous for initial submission).
\usepackage[accepted]{icml2026}

\icmltitlerunning{Activation Steering for Concept Unlearning in Diffusion Transformers}

\begin{document}

\twocolumn[
\icmltitle{Activation Steering for Concept Unlearning\\in Diffusion Transformers}

\icmlsetsymbol{equal}{*}

\begin{icmlauthorlist}
\icmlauthor{Anonymous Authors}{anon}
\end{icmlauthorlist}

\icmlaffiliation{anon}{Anonymous Institution, Anonymous City, Anonymous Country}
\icmlcorrespondingauthor{Anonymous Authors}{anon@example.com}

\icmlkeywords{Machine Learning, ICML, Mechanistic Interpretability,
Activation Steering, Concept Unlearning, Diffusion Transformers}

\vskip 0.3in
]

\printAffiliationsAndNotice{}

\begin{abstract}
Activation steering controls model behavior by adding a learned
direction to hidden states; it is well-studied in LLMs and recently
in U-Net diffusion models, but has not been characterized for modern
\emph{Diffusion Transformers} (DiTs). We present a training-free
concept-unlearning framework for DiTs that intervenes only at the
text-conditioning entry points: the CLIP-pooled embedding feeding
AdaLN modulation, and the projected T5 sequence entering joint
attention. A subspace zeroing diagnostic separates the two roles in
FLUX---CLIP gates object identity, T5 gates style---and motivates two
protocols. \emph{Style-Vacuum} subtracts a single time-averaged T5
direction. \emph{Object-Vacuum} subtracts a per-step T5 direction
together with an uncapped CLIP direction. With per-concept $\beta$
selection on UnlearnCanvas and TRACE-aligned LLaVA classification,
the method matches the published FLUX style-removal UA of TRACE
($88.9\%$ vs.\ $88.6\%$) and exceeds its FLUX object-removal UA
($97.3\%$ vs.\ $93.2\%$) without any retraining. The same recipe does
not transfer to Stable Diffusion~3.5; we report SD3.5 as a limitation
and discuss what its diagnostic shows about why.
\end{abstract}

\section{Introduction}
\label{sec:intro}

Text-to-image diffusion models have moved from U-Net backbones to
\emph{Diffusion Transformers}
(DiTs)~\cite{esser2024scaling,fluxdev2024}, which process text and
image tokens jointly across deep transformer stacks. Concept
unlearning---suppressing specific concepts in the output
distribution---is a standard safety/IP requirement, but most existing
methods retrain
weights~\cite{gandikota2023erasing,kumari2023ablating,gandikota2024unified,tracepaper}
or were developed for cross-attention U-Nets that do not survive joint
attention~\cite{khrulkov2024casteer}. We ask whether inference-time
activation steering---training-free, operating only on
activations---can do the same job in a DiT.

For FLUX, the answer is yes, and the reason is mechanistic. FLUX
exposes two text-conditioning entry points, and a simple diagnostic
shows each carries a distinct concept type: zeroing the CLIP-pooled
embedding on a fixed prompt collapses the generation to noise, while
zeroing the T5 path leaves the object intact and merely shifts style.
We design the unlearning procedure around this observation. Two
protocols are sufficient: \emph{Style-Vacuum} subtracts a single
T5 direction and \emph{Object-Vacuum} subtracts an uncapped CLIP
direction together with a per-step T5 direction. Per-concept $\beta$
selection on a small proxy set then handles the substantial
concept-to-concept variance that a single global $\beta$ leaves
unchecked. The same recipe does not transfer to SD3.5; we report this
as a limitation, with a diagnostic-grounded explanation in
\S\ref{sec:discuss}.

\paragraph{Contributions.}
(i)~Two diagnostic-grounded activation-steering protocols for FLUX
DiTs---\textbf{Style-Vacuum} and \textbf{Object-Vacuum}---intervening
only at text-conditioning entry points and requiring no model
retraining.
(ii)~A subspace zeroing diagnostic that grounds the protocol design
in where each concept type actually lives.
(iii)~A per-concept $\beta$ search via a 12-image proxy and a
symmetric UA/IRA/CRA composite, addressing the high variance of a
single global $\beta$ across $20$ objects and $10$ styles.
(iv)~Full UnlearnCanvas results on FLUX-schnell that match TRACE's
published FLUX style UA and exceed its object UA.

\section{Background and Related Work}
\label{sec:bg}

\paragraph{DiT architectures.}
FLUX and SD3.5 use \emph{Multi-Modal Diffusion Transformer} blocks in
which text and image tokens are concatenated and processed by joint
attention~\cite{esser2024scaling}. Both route text through two
parallel paths: a \emph{pooled} embedding that drives AdaLN
modulation across all blocks, and a \emph{sequence} embedding that
participates in joint attention with image tokens. FLUX uses a
$768$-d CLIP-L pooled vector and a T5 sequence, with $19$
double-stream and $38$ single-stream blocks; FLUX-schnell is a
$4$-step distillation that runs at guidance scale~$0$. SD3.5
concatenates two CLIP encoders into a $2048$-d pooled vector and
projects CLIP-L+G and T5 into a $333$-token sequence
($77$ CLIP $+\,256$ T5) processed by $24$ joint blocks under
classifier-free guidance.

\paragraph{Activation steering.}
In LLMs, activation steering adds a linear direction to hidden states
to control behavior~\cite{turner2024steering,panickssery2024steering};
the direction is the mean difference of activations between
contrastive prompts. CASteer~\cite{khrulkov2024casteer} extends the
idea to cross-attention U-Nets, introducing two ingredients we
adopt: $50$ ImageNet base contexts to diversify prompt pairs, and a
clip-negative ReLU rule that fires only when the activation is
positively aligned with the concept direction.

\paragraph{Concept unlearning.}
Most prior methods edit weights. ESD~\cite{gandikota2023erasing}
fine-tunes with negative prompts, UCE~\cite{gandikota2024unified}
performs closed-form weight edits, and TRACE~\cite{tracepaper}
replaces the projection layer and the modulation MLP with trained
transcoders, reporting the strongest UA on FLUX and SD3.5. TRACE's
App.~D.1 explicitly notes that, in SD3.5 and FLUX, ``despite operating
on a pooled representation,'' intervening at the modulation MLP is
necessary in addition to the projection layer for effective concept
suppression---a finding that motivates one of our diagnostics on
SD3.5. UnlearnCanvas~\cite{zhang2024unlearncanvas} provides the
$10$-style $\times\,20$-object grid we evaluate on, with three
metrics: \emph{Unlearning Accuracy} (UA, fraction of target-bearing
images not classified as the target), \emph{In-domain Retain
Accuracy} (IRA, accuracy on the target axis for non-target prompts),
and \emph{Cross-domain Retain Accuracy} (CRA, accuracy on the
orthogonal axis).

\section{Method}
\label{sec:method}

We treat concept unlearning as the problem of identifying and
neutralizing the text-conditioning \emph{entry points} through which
the concept reaches the joint-attention stack. The procedure has
four ingredients: a zeroing diagnostic that decides where to
intervene; a contrastive direction estimator at the chosen entry
points; the inference-time subtraction rule; and a per-concept
$\beta$ search on a small proxy. No model weights are touched.

\subsection{Subspace diagnostic}
\label{sec:method_diag}

For a fixed prompt $\langle\text{object},\text{style}\rangle$, we
generate the baseline once, then once per entry point with that
entry point's activation zeroed on the conditional batch position.
The panel reveals what each entry point carries.

For FLUX the result is unambiguous. Zeroing
\texttt{time\_text\_embed}'s input (the CLIP-pooled embedding)
collapses the generation to high-frequency noise; zeroing
\texttt{context\_embedder} (the T5 projection) preserves a
recognizable object and merely shifts the style. We read this as:
object identity is gated through CLIP\,$\to$\,AdaLN modulation; style
flows through T5\,$\to$\,joint attention. The two protocols below
follow directly.

\subsection{Direction estimation}
\label{sec:method_dirs}

For each concept $c$ we form $N{=}50$ contrastive prompt pairs using
$50$ ImageNet classes as base contexts, following
CASteer~\cite{khrulkov2024casteer}. For an object target,
$(p^+_i, p^-_i){=}(\text{``\texttt{class} with \texttt{concept}}'',\,
\text{``\texttt{class}}'')$; for a style target,
$(p^+_i, p^-_i){=}(\text{``\texttt{class}, \texttt{concept} style}'',\,
\text{``\texttt{class}}'')$. Averaging across diverse contexts cancels
prompt-specific signal and isolates the concept direction. The
direction at entry point $S$ is the L2-normalized mean of
$h_S(p^+_i)\!-\!h_S(p^-_i)$ across pairs.

\subsection{Two protocols}
\label{sec:method_protocols}

\paragraph{Style-Vacuum.} A single $4096$-d T5 direction is taken on
the \emph{input} of \texttt{context\_embedder} (raw T5 embeddings,
mask-aware mean-pooled across non-padding tokens). A $768$-d CLIP
direction on \texttt{time\_text\_embed} is also exposed, but the
diagnostic places style in T5 alone, so the default sets
$\beta_{\mathrm{CLIP}}{=}0$ and uses
$\beta_{\mathrm{T5}}\!\in\![1,6]$. Two vectors per concept; no
per-step factorization.

\paragraph{Object-Vacuum.} The CLIP direction is estimated as above.
For T5 we register a forward hook on the \emph{output} of
\texttt{context\_embedder} and accumulate per-step, mask-aware
mean-pooled differences over the denoising trajectory, yielding one
$3072$-d direction per step. A per-step direction is necessary
because object identity is committed within FLUX-schnell's $4$
denoising steps; a single time-averaged direction under-removes it.
CLIP is intentionally \emph{uncapped}---the cap parameter $\kappa$ in
Eq.~\ref{eq:steer} is $\infty$---so $\beta_{\mathrm{CLIP}}\!>\!1$ can
push the projection past zero into anti-concept space. Capping at
$\kappa{=}1$ subtracts only one axis of pooled CLIP, and because
object identity occupies a multi-dimensional CLIP subspace the model
merely rolls to a neighbouring class.

\subsection{Inference-time injection}
\label{sec:method_inject}

For activation $h_S$, direction $\hat v_S$, and strength
$\beta_S\!\geq\!0$, we apply CASteer's clip-negative rule:
\begin{equation}
\label{eq:steer}
h'_S \;=\; h_S \,-\, \min(\beta_S, \kappa)\cdot\mathrm{ReLU}\!\bigl(\langle h_S, \hat v_S\rangle\bigr) \cdot \hat v_S.
\end{equation}
The $\mathrm{ReLU}$ clamp ensures we only fire on positively aligned
activations, avoiding sign-flip amplification on unrelated prompts.
The cap $\kappa$ is per-protocol: $\kappa{=}1$ for Style-Vacuum
($\beta_{\mathrm{CLIP}}{=}0$ keeps the CLIP path inert anyway) and
$\kappa{=}\infty$ for Object-Vacuum. For sequence-valued entry points
(T5), the projection is computed per token; an optional top-fraction
gate localizes the subtraction to concept-aligned tokens. Steering
applies \emph{only} to the conditional batch position; the
unconditional CFG reference frame is left untouched. FLUX-schnell
runs at guidance scale~$0$, so the conditional position is the entire
forward.

\subsection{Per-concept $\beta$ selection}
\label{sec:method_hp}

A single global $\beta$ is empirically suboptimal. At
$(\beta_{\mathrm{CLIP}},\beta_{\mathrm{T5}}){=}(3,5)$ on FLUX,
distinctive silhouettes (Architectures, Bears) over-steer to
UA${=}100\%$ but CRA${\sim}30\%$, while broad concepts (Butterfly,
Flame) under-steer to UA${<}5\%$. We therefore run a per-concept
grid search on a $12$-image proxy ($3$ proxy styles
$\times\,2$ proxy objects $\times\,2$ seeds, with the target on the
relevant axis), classify proxy outputs with LLaVA, and score by
\[
\mathrm{score}(\beta) \;=\; 0.5\cdot\mathrm{UA} + 0.25\cdot\mathrm{IRA} + 0.25\cdot\mathrm{CRA}.
\]
The weighting is symmetric: UA dominant because it is the headline
metric, IRA and CRA equal because UnlearnCanvas explicitly shows the
two retention metrics dissociate. The object grid sweeps
$(\beta_{\mathrm{CLIP}},\beta_{\mathrm{T5}})\!\in\!
\{(1{,}3),(3{,}5),(5{,}8),(8{,}12),(10{,}15)\}$; the style grid
sweeps $\beta_{\mathrm{T5}}\!\in\!\{1,2,4,6\}$ at
$\beta_{\mathrm{CLIP}}{=}0$. Selected $\beta$'s are cached and
reused for the full benchmark.

\section{Experiments}
\label{sec:exp}

\subsection{Setup}
\label{sec:exp_setup}

\paragraph{Model.} All reported numbers are on FLUX.1-schnell
($4$ inference steps, guidance scale~$0$, license-permissive). Vector
estimation runs once per concept in ${\sim}100$ forward passes
($N{=}50$ pairs $\times\,2$ prompts) and the result is cached.

\paragraph{Benchmark and prompt.} We use the UnlearnCanvas
$10\!\times\!20$ grid restricted to the TRACE concept set, with the
prompt template ``A \texttt{\{object\}} image in
\texttt{\{style\}} style.''.

\paragraph{Classifier.} Generations are classified by
LLaVA-1.6-Vicuna-7B using the numbered-list prompt format introduced
by TRACE (App.~E.4); the original UnlearnCanvas SD-1.5 classifier
generalizes below $6\%$ to FLUX outputs and is unsuitable.

\paragraph{Seeds and grid size.} We evaluate on the five TRACE
reproducibility seeds $\{188, 288, 588, 688, 888\}$, generating
$10\!\times\!20\!\times\!5{=}1000$ images per target concept. We
additionally report FID against a shared $1000$-image unsteered
baseline grid (generated once and reused across concepts) and the
CLIP-ViT-L/14 score on the generation prompt.

\subsection{Object unlearning}
\label{sec:exp_obj}

Object-Vacuum yields an average UA of $\mathbf{97.3\%}$ across the
$20$ object concepts (Table~\ref{tab:obj}), with average
IRA${=}88.45\%$ (other-object preservation) and average
CRA${=}33.82\%$ (style preservation on non-target images). Fifteen of
twenty concepts reach UA${=}100\%$; the lowest UA values are on
Human ($86\%$) and Sea ($86\%$), categories where LLaVA confounds
with adjacent classes. Per-concept FID ranges from $31.6$ to $64.4$
and tracks the prompt-level disruption introduced by removing
scene-anchored objects (Sea and Sandwiches: low FID; Frogs and
Waterfalls: high FID). The headline UA exceeds TRACE's published
FLUX object UA ($93.2\%$); CRA is markedly lower
(\S\ref{sec:exp_compare}), an explicit trade-off we discuss in
\S\ref{sec:discuss}.

\begin{table}[t]
\caption{FLUX-schnell, object unlearning (Object-Vacuum), full
UnlearnCanvas grid. UA, IRA, CRA in $\%$; FID lower is better; CLIP
higher is better.}
\label{tab:obj}
\centering
\small
\begin{tabular}{lrrrrr}
\toprule
Concept & UA & IRA & CRA & FID & CLIP \\
\midrule
Architectures & 100  & 92.95 & 30.11 & 41.40 & 0.266 \\
Bears         & 100  & 88.63 & 35.37 & 42.28 & 0.263 \\
Birds         & 100  & 86.11 & 41.16 & 50.35 & 0.260 \\
Butterfly     & 100  & 89.26 & 32.11 & 40.92 & 0.261 \\
Cats          & 100  & 85.58 & 36.84 & 47.48 & 0.262 \\
Dogs          & 98   & 88.95 & 36.53 & 39.32 & 0.267 \\
Fishes        & 100  & 89.16 & 34.21 & 47.83 & 0.259 \\
Flame         & 98   & 88.32 & 27.68 & 41.54 & 0.263 \\
Flowers       & 96   & 87.89 & 30.42 & 51.20 & 0.260 \\
Frogs         & 100  & 85.05 & 35.16 & 64.42 & 0.234 \\
Horses        & 100  & 86.21 & 38.95 & 54.37 & 0.261 \\
Human         & 86   & 90.42 & 26.32 & 46.09 & 0.258 \\
Jellyfish     & 98   & 90.11 & 34.74 & 34.31 & 0.263 \\
Rabbits       & 100  & 88.74 & 37.58 & 53.38 & 0.250 \\
Sandwiches    & 100  & 87.37 & 38.95 & 44.03 & 0.265 \\
Sea           & 86   & 92.42 & 31.89 & 31.57 & 0.270 \\
Statues       & 92   & 88.84 & 34.74 & 38.72 & 0.268 \\
Towers        & 100  & 87.68 & 32.21 & 48.15 & 0.260 \\
Trees         & 100  & 89.58 & 34.21 & 44.77 & 0.263 \\
Waterfalls    & 100  & 85.79 & 27.26 & 56.89 & 0.253 \\
\midrule
\textbf{Average} & \textbf{97.30} & \textbf{88.45} & \textbf{33.82} & \textbf{45.95} & \textbf{0.260} \\
\bottomrule
\end{tabular}
\end{table}

\subsection{Style unlearning}
\label{sec:exp_style}

Style-Vacuum yields average UA${=}\mathbf{88.9\%}$ across the ten
TRACE-aligned styles (Table~\ref{tab:style}), statistically on par
with TRACE's published FLUX style UA of $88.6\%$, while leaving CRA
high ($86.6\%$). IRA at $35.72\%$ sits at the same structural ceiling
TRACE reports ($36.1\%$): when one style is removed and LLaVA is
forced to choose another from a $10$-way list, accuracy on the
remaining nine is bounded by the classifier's discrimination on
neutralized images. Cubism, Winter, Pop\_Art, and Ukiyoe each reach
UA${=}100\%$; Impressionism is the weakest (UA${=}66\%$), consistent
with its visual proximity to a generic realistic palette.

\begin{table}[t]
\caption{FLUX-schnell, style unlearning (Style-Vacuum), full
UnlearnCanvas grid.}
\label{tab:style}
\centering
\small
\begin{tabular}{lrrrrr}
\toprule
Concept & UA & IRA & CRA & FID & CLIP \\
\midrule
Van\_Gogh       & 86  & 38.89 & 88.78 & 31.56 & 0.267 \\
Watercolor      & 99  & 35.67 & 85.89 & 58.31 & 0.242 \\
Cartoon         & 85  & 30.67 & 87.33 & 65.11 & 0.240 \\
Cubism          & 100 & 40.22 & 84.11 & 75.21 & 0.221 \\
Winter          & 100 & 24.11 & 83.11 & 71.23 & 0.224 \\
Pop\_Art        & 100 & 38.56 & 87.67 & 59.70 & 0.239 \\
Ukiyoe          & 100 & 32.22 & 82.89 & 73.66 & 0.220 \\
Impressionism   & 66  & 35.67 & 87.11 & 39.10 & 0.262 \\
Byzantine       & 81  & 39.78 & 89.67 & 29.69 & 0.271 \\
Bricks          & 72  & 41.44 & 89.44 & 34.54 & 0.272 \\
\midrule
\textbf{Average} & \textbf{88.90} & \textbf{35.72} & \textbf{86.60} & \textbf{53.81} & \textbf{0.246} \\
\bottomrule
\end{tabular}
\end{table}

\subsection{Comparison with TRACE on FLUX}
\label{sec:exp_compare}

Table~\ref{tab:compare} reproduces the FLUX rows of TRACE's
Table~$1$~\cite{tracepaper} alongside our two protocols. On UA we
match TRACE on style and exceed it on objects, with no weight
updates. The remaining gap is on retention---particularly object
CRA, where weight-editing methods preserve more of the orthogonal
axis. We do not claim a Pareto improvement; the headline finding is
that an inference-time entry-point intervention reaches
transcoder-level UA on FLUX without any model fitting.

\begin{table}[t]
\caption{FLUX comparison against TRACE Table~$1$. Higher is better
for UA/IRA/CRA, lower for FID. ``--'' = not reported by TRACE.}
\label{tab:compare}
\centering
\small
\begin{tabular}{llrrrr}
\toprule
Task & Method & UA & IRA & CRA & FID \\
\midrule
\multirow{4}{*}{Style}
 & LOCOEDIT~\cite{tracepaper}      & 66.45 & 33.23 & 83.44 & 55.56 \\
 & UCE~\cite{gandikota2024unified} & 67.43 & 34.78 & 76.56 & 58.90 \\
 & TRACE~\cite{tracepaper}         & 88.60 & 36.10 & 96.40 & 51.67 \\
 & Style-Vacuum (ours)             & 88.90 & 35.72 & 86.60 & 53.81 \\
\midrule
\multirow{2}{*}{Object}
 & TRACE~\cite{tracepaper}         & 93.20 & 93.20 & 96.60 & --- \\
 & Object-Vacuum (ours)            & 97.30 & 88.45 & 33.82 & 45.95 \\
\bottomrule
\end{tabular}
\end{table}

\subsection{Ablations}
\label{sec:exp_abl}

\paragraph{Number of contrastive pairs.} Direction quality is
visibly noisier at $N{=}10$ and stabilizes by $N{=}50$, consistent
with CASteer's reported saturation~\cite{khrulkov2024casteer}.

\paragraph{CLIP cap.} Capping $\beta_{\mathrm{CLIP}}$ at
$\kappa{=}1$ silently throttles object steering: subtracting
$1\!\times\!\hat v_{\mathrm{CLIP}}$ zeroes one axis of pooled CLIP,
but object identity occupies a multi-dimensional subspace, so the
model rolls to a neighbouring class. Removing the cap and using
$\beta_{\mathrm{CLIP}}{=}3$ pushes the projection through zero and is
what produces clean object removal.

\paragraph{Composite objective.} Replacing the symmetric
$0.5{:}0.25{:}0.25$ composite with UA-only over-selects
retention-killing configurations on objects with distinctive
silhouettes; UA${+}$IRA under-selects on concepts where CRA is the
more pressing failure mode. The symmetric composite produces
selections closest to the visual sweet spot in our review.

\paragraph{Boundary diagnostics.} On roughly $15\%$ of FLUX object
concepts (Flame, Butterfly, Flowers), the per-concept search lands
on the boundary of the $5$-cell grid, indicating the true optimum
lies outside the searched range. We log these as ``boundary-warning''
picks; extending the grid resolves them at the cost of additional
search compute.

\section{Discussion and Limitations}
\label{sec:discuss}

The FLUX results carry a mechanistic reading rather than a purely
engineering one. The diagnostic isolates two text pathways with
disjoint roles---CLIP\,$\to$\,AdaLN for object identity,
T5\,$\to$\,joint attention for style---and the two-protocol design
follows directly. FLUX-schnell's distilled inference (no CFG
amplification) and double-stream attention let entry-point
subtractions propagate through AdaLN with minimal scrambling, which
is why training-free entry-point steering is competitive with
transcoder-based weight editing on UA.

\paragraph{SD3.5.} The same recipe does not transfer to
SD3.5-medium, and we do not report numerical results on it. Our
diagnostic shows two structural reasons. First, SD3.5
\emph{redundantly} encodes object identity across encoders: on the
prompt ``a Dog in Van Gogh style'', zeroing the full $2048$-d pooled
vector changes $96\%$ of pixels but leaves a recognizable animal;
the CLIP-L slice alone changes $36\%$, the CLIP-G slice $77\%$, and
in both cases object identity persists. There is no single
entry-point analog of FLUX's CLIP-pooled lever. Second, the
joint-attention stack appears to reconstruct concept features from
signal that survives input-side subtraction: at low and moderate
strengths, joint multi-subspace steering across the four input-side
subspaces (\texttt{pooled\_clipL}, \texttt{pooled\_clipG},
\texttt{ctx\_clip}, \texttt{ctx\_t5}) leaves the dog visible; at high
strengths the image structure collapses but residual object features
remain embedded in the new structure. Following TRACE App.~D.1, we
additionally implemented a hook on \texttt{tte\_out}---the post-SiLU
output of \texttt{time\_text\_embed}, i.e., the AdaLN modulation
signal that feeds every block---and the diagnostic confirms it acts
on the right surface, but tuning a competitive recipe was not
feasible within the project window. Closing this gap is the natural
next step.

\paragraph{Other limitations.} We evaluate on FLUX-schnell ($4$
steps, guidance~$0$); the absence of CFG and the lower step count
simplify entry-point steering relative to FLUX-dev ($28$ steps), and
absolute UA comparisons against TRACE's published FLUX row should be
read with this caveat. Object CRA is meaningfully lower than
transcoder-based methods because pushing the CLIP-pooled projection
past zero scrambles the AdaLN-modulated style signal that shares the
path. Finally, the contrastive-pair direction is rank-$1$ by
construction; distributed concepts may benefit from rank-$k$
subspace removal as in CASteer~\cite{khrulkov2024casteer}, which we
did not evaluate.

\section*{Software and Data}

Code, vector caches, and per-concept results CSVs will be released on
acceptance.

\bibliographystyle{icml2026}
\bibliography{references}

\appendix

\section{Implementation details}
\label{app:impl}

We use \texttt{FLUX.1-schnell} ($4$ inference steps, guidance
scale~$0$). All steering hooks operate exclusively on the
conditional batch position. LLaVA classifies into the $10$ TRACE
styles and the $20$ TRACE object classes with the numbered-list
prompt format. FID is computed against a shared $1000$-image
unsteered baseline grid produced once and reused across all
concepts. The full reproducibility configuration, including the
per-concept $\beta$ JSON cache, is in the code repository.

\end{document}
